\scanpage{115}
Soit $A$ un anneau commutatif, $X$ le schéma affine qu'il définit,
\[
A^{\mathrm{cons}}=\underset{\unsure{\mathfrak P|\mathcal P}}{\unsure{\underline{\lim}|\varinjlim}} A_{\unsure{\mathfrak P|\mathcal P}}
\]
$\unsure{\mathfrak P|\mathcal P}$ parcourt les ``partitions \unsure{ordinaires} de $X$ en ens constructibles loc. fermés, affines, et pour une telle $\unsure{\mathfrak P|\mathcal P}=(X_i)_{i\in I}$, $A_{\unsure{\mathfrak P|\mathcal P}}$ désigne \deleted{\unsure{le produit}} $A(\coprod X_{i\,\mathrm{red}})=\prod A(X_{i\,\mathrm{red}})$.

Les \unsure{hom.} de transition sont injectifs. Considérons
\[
\widetilde X=\operatorname{Spec}(A^{\mathrm{cons}})
\]
On a
\[
\operatorname{Top}(\widetilde X)=\mathop{\underleftarrow{\mathrm{Lim}}}_{\unsure{\mathfrak P|\mathcal P}}\operatorname{Top}\operatorname{Sp}(A_{\unsure{\mathfrak P|\mathcal P}})
\]
En particulier
\[
\operatorname{Ens}(\widetilde X)=\mathop{\underleftarrow{\mathrm{Lim}}}_{\unsure{\mathfrak P|\mathcal P}}\underbrace{\operatorname{Ens}\operatorname{Sp}A_{\unsure{\mathfrak P|\mathcal P}}}_{\coprod X_i}
\]
et comme les morphismes de transition sont des \deleted{\unsure{isomorphismes}} \inserted{bijections}, on trouve \deleted{\unsure{par}} une \deleted{\unsure{surj}} \inserted{bijection}.
\[
\operatorname{Ens}\widetilde X\simeq\operatorname{Ens}X
\]
Cette bijection est compatible avec les topologies de $\widetilde X$ et celle de $X^{\mathrm{cons}}$, comme on voit aisément sur la description d'une base d'ouverts de $\widetilde X$ ....
\margin{114}
